% =============================================================================
%  CV de Emanuel Edgar Ancco Guaygua - Beca AECODE 2026
%  Enfoque: Investigación, Desarrollo Tecnológico e Innovación (I+D+i)
%  Versión: Final AECODE (Optimized Skills & Project Details)
% =============================================================================

\documentclass[10pt, a4paper]{article}

% --- PAQUETES NECESARIOS ---
\usepackage[utf8]{inputenc} 
\usepackage[spanish]{babel}      
\usepackage{geometry}            
\usepackage{hyperref}            
\usepackage{fontawesome5}        
\usepackage{titlesec}            
\usepackage{amsmath}             
\usepackage{lmodern} 

% --- CONFIGURACIÓN DE LA PÁGINA ---
\geometry{
    a4paper,
    left=1.6cm,
    right=1.6cm,
    top=1.3cm,
    bottom=1.3cm
}

% --- COLORES Y ENLACES ---
\hypersetup{
    colorlinks=true,
    linkcolor=black,
    urlcolor=blue,
    pdftitle={CV - Emanuel Edgar Ancco Guaygua - AECODE},
    pdfauthor={Emanuel Edgar Ancco Guaygua},
}

% --- PERSONALIZACIÓN DE TÍTULOS DE SECCIÓN ---
\titleformat{\section}
  {\Large\bfseries\scshape} 
  {}
  {0em}
  {}[\titlerule] 

\titlespacing*{\section}{0pt}{2.0ex plus 0.8ex minus .3ex}{1.2ex plus .3ex}

% --- ELIMINAR SANGRÍA DE PÁRRAFO ---
\setlength{\parindent}{0pt}

% --- DEFINICIÓN DE COMANDOS PARA ENTRADAS ---
\newcommand{\entry}[4]{
    \textbf{#1} \hfill \textit{#2} \\
    \textit{#3} \hfill \textit{#4} \\
    \vspace{-1.5mm} 
}

% =============================================================================
%  INICIO DEL DOCUMENTO
% =============================================================================
\begin{document}

% --- ENCABEZADO ---
\begin{center}
    {\Huge\bfseries EMANUEL EDGAR ANCCO GUAYGUA}
    \vspace{2mm}

    \faMapMarker*{} Lima / Puno, Perú \quad
    \faPhone*{} 974583549 \quad
    \faEnvelope{} \href{mailto:coarp.eancco@gmail.com}{coarp.eancco@gmail.com} \quad
    \faLinkedin{} \href{https://www.linkedin.com/in/emanuel-edgar-ancco-guaygua-a61210352}{LinkedIn}
    
    \faGlobe{} \textbf{Portafolio de Investigación:} \href{https://emanuelancco.github.io/cv-portfolio/}{emanuelancco.github.io/cv-portfolio}
\end{center}

\vspace{2mm}

% --- RESUMEN PROFESIONAL ---
\section*{Perfil de Investigador y Desarrollador}
Ingeniero Civil (UPC, Diciembre 2025) especializado en \textbf{Inteligencia Artificial aplicada a Infraestructura Crítica} y desarrollo de software de ingeniería. Líder en investigación de sistemas de monitoreo estructural (SHM) utilizando \textbf{Deep Learning (GNN, PINN)}, Visión Artificial avanzada y computación en el borde (Edge AI). Desarrollador full-stack con dominio de \textbf{Python, C\# y Arquitecturas de Hardware Reconfigurable (FPGA)}. Experto en gestión de proyectos complejos (\textbf{Primavera P6, Delphin Express}) y metodologías BIM/VDC. Ganador de concursos de innovación internacional (DigiEduHack) y ponente académico (CONEIC/COREIC).

\vspace{2mm}
\begin{center}
\fbox{\parbox{0.92\textwidth}{\centering
\faExternalLink*{} \textbf{EVIDENCIA TÉCNICA Y DEMOS:} \href{https://emanuelancco.github.io/cv-portfolio/}{emanuelancco.github.io/cv-portfolio} \\
\small{Videos de inferencia, métricas de modelos ML, arquitecturas GNN, capturas de plugins Revit}
}}
\end{center}

\vspace{3mm}

% --- EXPERIENCIA PREPROFESIONAL ---
\section*{Experiencia Profesional y de Investigación}

\entry{Asistente de Oficina Técnica - I+D}{Septiembre 2025 – Diciembre 2025}{Obra: ``Creación de Servicios Culturales - Teatro Municipal de Yunguyo''}{Puno, Perú}
\begin{itemize}
    \item \textbf{Automatización BIM:} Desarrollo de plugins personalizados en \textbf{Revit API 2026 (C\#)} para generación automática de metrados y armaduras estructurales.
    \item \textbf{Gestión de Proyectos:} Control de plazos y costos implementando \textbf{Primavera P6, Microsoft Project y Delphin Express}. Optimización del flujo de valorizaciones mediante scripts de Python.
    \item \textbf{Ingeniería de Costos:} Elaboración de expedientes técnicos y presupuestos utilizando \textbf{SAP (Nivel Intermedio)} y S10.
\end{itemize}

\vspace{2mm} 

\entry{Asistente de Oficina Técnica}{Junio 2024 – Noviembre 2024}{Constructora Talgus S.A.C.}{Lima, Perú}
\begin{itemize}
    \item Elaboración de metrados y análisis de costos unitarios para preparación de presupuestos.
    \item Seguimiento y control de costos de obra, preparación de informes de valorización.
    \item Modelado y actualización de planos técnicos (AutoCAD, Revit) para correcta ejecución.
\end{itemize}

\vspace{2mm} 

\entry{Director del Área de Investigaciones}{Enero 2024 – Diciembre 2025}{Capítulo Estudiantil de Ingeniería Civil - IDI UPC}{Lima, Perú}
\begin{itemize}
    \item \textbf{Liderazgo Técnico:} Dirección de equipos multidisciplinarios para desarrollo de soluciones IA en ingeniería civil (Detección de patologías, SHM).
    \item \textbf{Investigación Aplicada:} Supervisión del ciclo completo de proyectos I+D, desde validación de hipótesis hasta implementación de prototipos (Simuladores, Modelos ML).
    \item \textbf{Difusión Académica:} Ponente en CONEIC Ucayali 2024 sobre ``Aplicación de CNN para Evaluación Automatizada de Fisuras y Grietas en Construcciones de Lima''.
\end{itemize}

\vspace{3mm}

% --- EDUCACIÓN ---
\section*{Educación}
\entry{Pregrado en Ingeniería Civil (Egresado)}{2021 - Diciembre 2025}{Universidad Peruana de Ciencias Aplicadas (UPC)}{Lima, Perú}
\vspace{1mm}
\entry{Bachillerato Internacional (IB Diploma)}{2015 - 2017}{Colegio de Alto Rendimiento (COAR Puno)}{Puno, Perú}

\vspace{3mm}

% --- PROYECTOS DE DESARROLLO DE SOFTWARE E IA ---
\section*{Proyectos de Desarrollo Tecnológico e IA}

\textbf{GAIATECH SHM | Sistema de Monitoreo de Salud Estructural con IA}
\begin{itemize}
    \item Sistema avanzado de \textbf{detección de anomalías en puentes} usando datos de 5 acelerómetros y redes neuronales de grafos.
    \item \textbf{Innovación:} Arquitectura \textbf{PINN-GNN} (Physics-Informed Graph Neural Network) que integra leyes físicas en el grafo de aprendizaje.
    \item Mejor modelo: \textbf{STG-AE Wavelet-GNN} con 5.1M parámetros y \textit{validation loss} de \textbf{0.0064}.
    \item Extracción de features con \textbf{Discrete Wavelet Transform (DWT)} y análisis espectral (FFT, CWT).
    \item Tecnologías: Python, \textbf{PyTorch, PyTorch Geometric}, Wavelets (PyWavelets), CUDA.
\end{itemize}

\vspace{2mm}

\textbf{EMARC VISIÓN | Sistema de Detección de Grietas y EPP con Visión Artificial}
\begin{itemize}
    \item \textbf{Módulo de Grietas (SegFormer):} Segmentación semántica de alta precisión para fisuras en concreto. Clasificación automática de severidad (Leve/Moderada/Severa).
    \item Cálculo automático de métricas físicas: \textbf{área (cm²), longitud, ancho promedio, porcentaje de daño}.
    \item Recomendaciones automáticas de mantenimiento según normativa (inspección/reparación).
    \item \textbf{Módulo de Seguridad (YOLOv8):} Detección de EPP en tiempo real con \textbf{mAP@0.5 = 90.7\%}.
    \item Tecnologías: Python, \textbf{Hugging Face Transformers}, SegFormer, YOLOv8, OpenCV, Tkinter GUI.
\end{itemize}

\vspace{2mm}

\textbf{Suite de Análisis de Tuberías Enterradas (Ondas TGD) | Software de Ingeniería}
\begin{itemize}
    \item \textbf{v1 - Interfaz de Diseño Avanzada (Python/Tkinter):} Software de escritorio para \textbf{análisis de esfuerzos por tramos}. Incorpora visualización interactiva de distribución de esfuerzos mediante \textbf{Gráficos Polares y Modelado 3D} para identificación de puntos críticos en la red.
    \item \textbf{v2 - Módulo de Investigación Probabilística:} Simulador estocástico para evaluación de confiabilidad sísmica. Implementación de \textbf{Simulación Monte Carlo y método FORM}. Generación de \textbf{Curvas de Fragilidad} y análisis de corrosión dependiente del tiempo.
    \item Tecnologías: Python, SciPy (Stats), Pandas, Matplotlib (Polar/3D), Análisis de Riesgo Sísmico.
\end{itemize}

\vspace{2mm}

\textbf{Plataforma de Predicción de Riesgos Laborales | NLP + Deep Learning}
\begin{itemize}
    \item API REST (FastAPI) para clasificación de narrativas de incidentes usando \textbf{DistilBERT}.
    \item Predicción de naturaleza del evento, parte del cuerpo afectada y tipo de evento.
    \item Módulo de \textbf{Toolbox Talks} automáticos para prevención proactiva de accidentes.
    \item Tecnologías: Python, \textbf{PyTorch, Transformers (Hugging Face)}, FastAPI, SQL.
\end{itemize}

\vspace{2mm}

\textbf{Desarrollo de Asistentes IA y Agrotecnología}
\begin{itemize}
    \item \textbf{Archon Assistant:} Asistente virtual de voz desarrollado en Python. Integración de \textbf{OpenAI API (GPT)} para razonamiento complejo, librerías de reconocimiento de voz (\textit{SpeechRecognition}) y síntesis de voz (\textit{pyttsx3}).
    \item \textbf{AgroTech Vision:} Sistema de detección temprana de enfermedades en cultivos (Papa) utilizando \textbf{Redes Neuronales Convolucionales (CNN)}. Entrenado con dataset PlantVillage para clasificación de tizón tardío/temprano.
\end{itemize}

\vspace{2mm}

\textbf{EMARC BIM SUITE | Suite de Automatización para Revit 2026}
\begin{itemize}
    \item Suite de plugins de alto rendimiento desarrollados en \textbf{Visual Studio (C\#, WPF)}.
    \item \textbf{SAPtoRevit:} Importador de modelos SAP2000/ETABS a Revit (vigas, columnas, losas con pendientes).
    \item \textbf{EMARC Structural:} Suite de 20+ comandos para armado automático de elementos estructurales.
    \item \textbf{AutoCAD $\leftrightarrow$ Revit:} Interoperabilidad JSON para transferencia de geometría y familias.
    \item Tecnologías: C\#, \textbf{Revit API 2026}, .NET Framework 4.8, Visual Studio 2022.
\end{itemize}

\vspace{2mm}

\textbf{EMAIRC HIDRA | Suite de Ingeniería Hidráulica para HP Prime}
\begin{itemize}
    \item Suite de cálculo hidráulico programada en \textbf{PPL (Pascal)} con arquitectura de loop no recursivo y librería GUI personalizada táctil.
    \item Implementa métodos numéricos iterativos (Secante, Standard Step) para cálculo de curvas de remanso, tirantes conjugados y flujo variado con integración gráfica.
\end{itemize}

\vspace{2mm}

\textbf{Innovación Tecnológica e Instrumentación IoT (Hardware Propio)}
\begin{itemize}
    \item \textbf{PachaGuard (IoT/FPGA):} Sistema de vigilancia autónoma con \textbf{FPGA Tang Nano 1K} y ESP32-CAM. Procesamiento en el borde (Edge Computing) para alertas en tiempo real vía UART. Código en Verilog.
    \item \textbf{Detección Ambiental (Aves/Tráfico):} Modelos \textbf{YOLOv11 y DETR} (Transformers) para monitoreo de fauna y flujo vehicular.
    \item \textbf{Análisis de Materiales:} Scripting avanzado con librería \textit{concrete-properties} (ACI) para diagramas momento-curvatura.
    \item Tecnologías: Verilog (FPGA), C++ (Arduino), MicroPython.
\end{itemize}

\vspace{2mm}

\textbf{Herramientas de Análisis Estructural e Hidráulico}
\begin{itemize}
    \item \textbf{Análisis Sísmico Modal:} Herramienta Python para análisis dinámico con norma E.030 (matrices de rigidez, modos, derivas).
    \item \textbf{Simulación Monte Carlo:} Análisis de confiabilidad estructural con métodos FORM/FOSM y curvas de fragilidad.
    \item \textbf{Diseño Hidráulico:} Cálculo de canales con ecuación de Manning, exportación a Civil 3D, HEC-RAS, IBER.
    \item \textbf{Detección de Tráfico Vehicular:} Sistema YOLOv5/v11 para conteo y clasificación de vehículos en tiempo real.
    \item Tecnologías: Python, NumPy, SciPy, PyWavelets, OpenCV.
\end{itemize}

\vspace{3mm}

% --- HABILIDADES TÉCNICAS ---
\section*{Habilidades Técnicas}
\begin{itemize}
    \item \textbf{Lenguajes de Programación:} Python (Experto), C\# (.NET), C++, Verilog, PPL (HP Prime), SQL.
    \item \textbf{Inteligencia Artificial:} PyTorch, TensorFlow, Scikit-learn, YOLOv8/v11, SegFormer, Hugging Face Transformers.
    \item \textbf{Ingeniería BIM/CAD:} Revit (Desarrollo de Plugins), AutoCAD, Civil 3D, SAP2000.
    \item \textbf{Gemelos Digitales & Web 3D:} Desarrollo de experiencias inmersivas con \textbf{Three.js/React} (Estilo Noomo), scripting en \textbf{Blender}.
    \item \textbf{Gestión de Proyectos:} \textbf{Primavera P6, MS Project, Delphin Express}, SAP (Nivel Intermedio).
    \item \textbf{Análisis de Datos:} R, MATLAB, SQL, Power BI, Excel VBA.
    \item \textbf{Confiabilidad Estructural:} Métodos FORM/FOSM, Simulación Monte Carlo, Curvas de Fragilidad.
    \item \textbf{Códigos de Diseño:} ACI 318, AISC 360, E.030, E.060, ASME B31.
    \item \textbf{Automatización:} n8n (workflows), Dynamo, generación automática de reportes Word/Excel.
    \item \textbf{Idiomas:} Español (Nativo), Inglés Intermedio-Avanzado (dominio técnico y conversacional).
\end{itemize}

\vspace{3mm}

% --- CERTIFICACIONES Y RECONOCIMIENTOS ---
\section*{Certificaciones y Reconocimientos}
\textit{\small Verificables en LinkedIn: \href{https://www.linkedin.com/in/emanuel-edgar-ancco-guaygua-a61210352}{linkedin.com/in/emanuel-edgar-ancco-guaygua}}
\vspace{2mm}
\begin{itemize}
    \item \textbf{Especialización en Machine Learning on Google Cloud} (4 cursos) | Google Cloud / Coursera, 2025
    \item \textbf{Especialista en IA y Aprendizaje Automático con Python} (120 horas) | Data Science Analysis, 2024
    \item \textbf{Ponente en CONEIC Ucayali 2024} - ``Aplicación de CNN para Evaluación Automatizada de Fisuras y Grietas en Construcciones de Lima''
    \item \textbf{2do lugar en Concurso de Conocimientos CONEIC} | Huaraz, 2025
    \item \textbf{Primer Puesto, Concurso de Ponencias Académicas} | XIV COREIC Lima 2023, UNFV
    \item \textbf{Innovador Destacado, DigiEduHack 2023} | Comisión Europea
    \item \textbf{Especialización en Python for Everybody} (5 cursos) | University of Michigan / Coursera
    \item \textbf{ISC AI Foundations For Everyone} | Campus Global ISC, 2025
    \item \textbf{ISC Data Analytics Aplicado a Cadena de Suministro} | Campus Global ISC, 2025
\end{itemize}

\vspace{3mm}

% --- INFORMACIÓN ADICIONAL ---
\section*{Información Adicional}
\begin{itemize}
    \item \textbf{Disponibilidad:} Inmediata para programas de investigación o becas internacionales.
    \item \textbf{Origen:} Puno, Perú - \textbf{Familiarizado con trabajo en condiciones de altura} (4,000+ msnm).
    \item \textbf{Conocimiento normativo:} Familiaridad con estándares GISTM, CDA y ANCOLD para seguridad de presas.
\end{itemize}

\end{document}
